\chapter*{来源与证据}
\addcontentsline{toc}{chapter}{来源与证据}
公开模型事实优先引用官方模型卡、配置、代码和论文。模型系列公开报告不能替代特定检查点完整数据配方。本地参数、样本数量、失败处理与运行状态取自本项目源码和固定快照；建议方案与教学示例已在正文标明。网页访问日期为2026年11月14日。
\begingroup\renewcommand{\bibname}{公开一手资料}
\begin{thebibliography}{99}
\bibitem{qwen-card} Qwen/Qwen3-VL-2B-Instruct official model card.
\url{https://huggingface.co/Qwen/Qwen3-VL-2B-Instruct}
\bibitem{qwen-report} Qwen3-VL Technical Report, arXiv:2511.21631v2.
\url{https://arxiv.org/pdf/2511.21631v2}
\bibitem{qwen-repo} QwenLM/Qwen3-VL official repository.
\url{https://github.com/QwenLM/Qwen3-VL}
\bibitem{qwen-config} Qwen3-VL-2B-Instruct official model configuration.
\url{https://huggingface.co/Qwen/Qwen3-VL-2B-Instruct/blob/main/config.json}
\bibitem{multimodal-template} Hugging Face Transformers: Multimodal chat templates.
\url{https://huggingface.co/docs/transformers/chat_templating_multimodal}
\bibitem{qwen-transformers} Hugging Face Transformers: Qwen3-VL model documentation.
\url{https://huggingface.co/docs/transformers/model_doc/qwen3_vl}
\bibitem{lora-paper} Hu et al. — LoRA: Low-Rank Adaptation of Large Language Models (2021).
\url{https://arxiv.org/abs/2106.09685}
\bibitem{qlora-paper} Dettmers et al. — QLoRA: Efficient Finetuning of Quantized LLMs (2023).
\url{https://arxiv.org/abs/2305.14314}
\bibitem{peft-quantization} Hugging Face PEFT documentation — Quantization.
\url{https://huggingface.co/docs/peft/developer_guides/quantization}
\bibitem{transfer-feature-cache} François Chollet / Keras — Transfer learning \& fine-tuning.
\url{https://keras.io/guides/transfer_learning/}
\bibitem{dpo-paper} Rafailov et al. — Direct Preference Optimization: Your Language Model is Secretly a Reward Model (2023).
\url{https://arxiv.org/abs/2305.18290}
\bibitem{instructgpt-paper} Ouyang et al. — Training language models to follow instructions with human feedback (2022).
\url{https://arxiv.org/abs/2203.02155}
\end{thebibliography}\endgroup
\section*{本地证据}
\addcontentsline{toc}{section}{本地证据}
项目根目录：
\begin{Verbatim}
/home/luckysir/Strands_Robots
\end{Verbatim}
以下证据以项目根目录为起点：
\begin{itemize}
\item 三组采集计划、失败案例与原始资产审计：\path{experiments/microduck_tasks/outputs/training_collection_audit_20261114.json}
\item 正式数据划分、样本数量、标签分布与源记录：\path{experiments/microduck_tasks/outputs/training_dataset_final_20261114/manifest.json}
\item 数据泄漏、图片一致性与导出结构审计：\path{experiments/microduck_tasks/outputs/training_audit_20261114/final_dataset.json}
\item 模型、依赖、超参数、数据与权重指纹：\path{experiments/microduck_tasks/outputs/qlora_20261114/training_config.json}
\item 真实更新损失、学习率、显存与定期验证：\path{experiments/microduck_tasks/outputs/qlora_20261114/metrics.jsonl}
\item 正式首个检查点参数与状态核验：\path{experiments/microduck_tasks/outputs/qlora_20261114/startup_verification.json}
\item 资源试跑和恢复验证；与正式任务成绩分开：\path{experiments/microduck_tasks/outputs/resource_verification_20261114/resource_verification.json}
\end{itemize}
报告目录中的 \texttt{sources\_local\_data.json} 列出数据与执行代码的具体行号及SHA256；\texttt{sources\_methods.json} 列出训练方法来源；\texttt{snapshot\_*.json} 与 \texttt{snapshot\_metrics.jsonl} 是生成本书时使用的固定状态。原始训练日志在继续增长，报告不会自动刷新。
